% GENERATED FILE. Edit canonical lessons, then run npm run generate:books.
% canonical-source-hash: fnv1a64:8768e4a8f4dd1dd4
% canonical-lessons: SA-C142-tusyati, SA-C142-sramyati, SA-C142-racayati, SA-C142-darsayati, SA-C142-pujayati

\chapter{Is pleased, Gets tired, Composes, Shows, Worships}
\label{ch:sa-is-pleased-gets-tired-composes-shows-worships}
\hlchaptermodality{142}

\begin{chapteropening}
\textbf{By the end of this chapter:} \emph{I can say is pleased, gets tired, composes, shows and worships.}

Complete the last lesson of chapter 142: I can say is pleased, gets tired, composes, shows and worships.
\end{chapteropening}

\section[\texorpdfstring{\sk{तुष्यति} (tuṣyati)}{तुष्यति (tuṣyati)}]{\sk{तुष्यति} (tuṣyati) — is pleased, is satisfied}

\label{lesson:SA-C142-tusyati}



\begin{quote}
Before the new one: say the Sanskrit for touches, then the Sanskrit for leaves.
\end{quote}

\subsection*{You'll want to know: \sk{तुष्यति}}
\textbf{\sk{तुष्यति}} — \emph{tuṣyati} — \textquotedblleft{}is pleased, is satisfied\textquotedblright{}.

\begin{grammarlens}[title={What you've built}]
One more word for this chapter.
\end{grammarlens}

\subsection*{Your turn}
\begin{itemize}
\raggedright
  \item \emph{Say it:} \emph{tuṣyati}
  \item \emph{Say it:} \emph{tuṣyati}, once more
  \item \emph{Say it:} say \emph{tuṣyati}
  \item \emph{Recall:} say the Sanskrit for touches, then the Sanskrit for leaves, then say \emph{tuṣyati} again
\end{itemize}

\begin{tcolorbox}[breakable,colback=teal!4,colframe=teal!35!black,title={Before you move on}]
What does \sk{तुष्यति} mean? (\textquotedblleft{}Is pleased, is satisfied\textquotedblright{}.) Say it once more.
\end{tcolorbox}
\section[\texorpdfstring{\sk{श्राम्यति} (śrāmyati)}{श्राम्यति (śrāmyati)}]{\sk{श्राम्यति} (śrāmyati) — gets tired}

\label{lesson:SA-C142-sramyati}



\begin{quote}
Before the new one: say the Sanskrit for leaves, then the Sanskrit for is pleased.
\end{quote}

\subsection*{You'll want to know: \sk{श्राम्यति}}
\textbf{\sk{श्राम्यति}} — \emph{śrāmyati} — \textquotedblleft{}gets tired\textquotedblright{}.

\begin{grammarlens}[title={What you've built}]
One more word for this chapter.
\end{grammarlens}

\subsection*{Your turn}
\begin{itemize}
\raggedright
  \item \emph{Say it:} \emph{śrāmyati}
  \item \emph{Say it:} \emph{śrāmyati}, once more
  \item \emph{Say it:} say \emph{śrāmyati}
  \item \emph{Recall:} say the Sanskrit for leaves, then the Sanskrit for is pleased, then say \emph{śrāmyati} again
\end{itemize}

\begin{tcolorbox}[breakable,colback=teal!4,colframe=teal!35!black,title={Before you move on}]
What does \sk{श्राम्यति} mean? (\textquotedblleft{}Gets tired\textquotedblright{}.) Say it once more.
\end{tcolorbox}
\section[\texorpdfstring{\sk{रचयति} (racayati)}{रचयति (racayati)}]{\sk{रचयति} (racayati) — composes, makes}

\label{lesson:SA-C142-racayati}



\begin{quote}
Before the new one: say the Sanskrit for is pleased, then the Sanskrit for gets tired.
\end{quote}

\subsection*{You'll want to know: \sk{रचयति}}
\textbf{\sk{रचयति}} — \emph{racayati} — \textquotedblleft{}composes, makes\textquotedblright{}.

\begin{grammarlens}[title={What you've built}]
One more word for this chapter.
\end{grammarlens}

\subsection*{Your turn}
\begin{itemize}
\raggedright
  \item \emph{Say it:} \emph{racayati}
  \item \emph{Say it:} \emph{racayati}, once more
  \item \emph{Say it:} say \emph{racayati}
  \item \emph{Recall:} say the Sanskrit for is pleased, then the Sanskrit for gets tired, then say \emph{racayati} again
\end{itemize}

\begin{tcolorbox}[breakable,colback=teal!4,colframe=teal!35!black,title={Before you move on}]
What does \sk{रचयति} mean? (\textquotedblleft{}Composes, makes\textquotedblright{}.) Say it once more.
\end{tcolorbox}
\section[\texorpdfstring{\sk{दर्शयति} (darśayati)}{दर्शयति (darśayati)}]{\sk{दर्शयति} (darśayati) — shows}

\label{lesson:SA-C142-darsayati}



\begin{quote}
Before the new one: say the Sanskrit for gets tired, then the Sanskrit for composes.
\end{quote}

\subsection*{You'll want to know: \sk{दर्शयति}}
\textbf{\sk{दर्शयति}} — \emph{darśayati} — \textquotedblleft{}shows\textquotedblright{}.

\begin{grammarlens}[title={What you've built}]
One more word for this chapter.
\end{grammarlens}

\subsection*{Your turn}
\begin{itemize}
\raggedright
  \item \emph{Say it:} \emph{darśayati}
  \item \emph{Say it:} \emph{darśayati}, once more
  \item \emph{Say it:} say \emph{darśayati}
  \item \emph{Recall:} say the Sanskrit for gets tired, then the Sanskrit for composes, then say \emph{darśayati} again
\end{itemize}

\begin{tcolorbox}[breakable,colback=teal!4,colframe=teal!35!black,title={Before you move on}]
What does \sk{दर्शयति} mean? (\textquotedblleft{}Shows\textquotedblright{}.) Say it once more.
\end{tcolorbox}
\section[\texorpdfstring{\sk{पूजयति} (pūjayati)}{पूजयति (pūjayati)}]{\sk{पूजयति} (pūjayati) — worships}

\label{lesson:SA-C142-pujayati}



\begin{quote}
Before the new one: say the Sanskrit for composes, then the Sanskrit for shows.
\end{quote}

\subsection*{You'll want to know: \sk{पूजयति}}
\textbf{\sk{पूजयति}} — \emph{pūjayati} — \textquotedblleft{}worships\textquotedblright{}.

\begin{grammarlens}[title={What you've built}]
One more word for this chapter.
\end{grammarlens}

\subsection*{Your turn}
\begin{itemize}
\raggedright
  \item \emph{Say it:} \emph{pūjayati}
  \item \emph{Say it:} \emph{pūjayati}, once more
  \item \emph{Say it:} say \emph{pūjayati}
  \item \emph{Recall:} say the Sanskrit for composes, then the Sanskrit for shows, then say \emph{pūjayati} again
\end{itemize}

\begin{tcolorbox}[breakable,colback=teal!4,colframe=teal!35!black,title={Before you move on}]
What does \sk{पूजयति} mean? (\textquotedblleft{}Worships\textquotedblright{}.) Say it once more.
\end{tcolorbox}
